% Part: counterfactuals
% Chapter: introduction
% Section: paradoxes-material

\documentclass[../../../include/open-logic-section]{subfiles}

\begin{document}

\olfileid{cnt}{int}{par}

\olsection{د مادي شرطي پاراډوکسونه}

د انګريزۍ د «که \dots نو \dots» ويناوو د نښه‌ليکنې لپاره د
$!A \lif !B$ پر کارولو له لومړنيو نيوکه کوونکو څخه سي. آي.
لويس و۔ لويس د وايټ هېډ او رسل په \emph{پرنسيپيا متماتيکا}
کښې د مادي شرطي پر کارولو نيوکه کوله؛ هغوے $\lif$ د
«لازموي» په توګه لوسته۔ لويس په سمه توګه اعتراض وکړ چې که
$\lif$ د «لازموي» معنا ولري، نو هره کاذبه قضيه~$p$ دا
لازموي چې $p$،~$q$ لازموي، ځکه که~$p$ کاذب وي،
$p \lif (p \lif q)$ رښتيا دے؛ او هره رښتينې قضيه~$q$ دا
لازموي چې $p$،~$q$ لازموي، ځکه که $q$~رښتيا وي،
$q \lif (p \lif q)$ رښتيا دے۔

منطقپوهان البته پوهېږي چې \emph{لازمول}، يعنې منطقي
استلزام، نښلوونکے نۀ دے، بلکې د !!{formula}s يا ويناوو
ترمنځ اړيکه ده۔ نو د ګډوډۍ د مخنيوي لپاره بايد $\lif$ د
«لازموي» په توګه ونۀ لولو۔\footnote{رياضيپوهان او د
  کمپيوټرپوهې پوهان لا هم ډېر ځله «$\lif$» د «لازموي» په
  توګه لولي؛ خو فيلسوفان يې «يوازې که» بولي، څو د لويس له
  ښودلې ګډوډۍ څخه ډډه وکړي۔} تر څو چې داسې يې نۀ لولو، د
لويس ځانګړې اندېښنه نۀ راپيدا کېږي: $p$، $q$ نۀ «لازموي»،
که څه هم $p$ د يوې کاذبې انګريزي جملې نښه وګڼو۔ د دې معلومولو
لپاره چې آيا $p \Entails q$، بايد \emph{ټولې}
!!{valuation}s وګورو؛ او $p \Entails/ q$ دے، حتی هغه وخت هم
چې $p$ د داسې جملې د نښې په توګه کاروو چې په اتفاق کاذبه وي۔

خو بيا هم د لويس په شان د «که \dots نو\dots» ويناوو کښې يو
څه عجيب شته:
\begin{quote}
که سپوږمۍ له شنې پنيرې جوړه وي، نو $2+2=4$۔
\end{quote}
او په دغو استنتاجونو کښې هم:
\begin{quote}
  سپوږمۍ له شنې پنيرې جوړه نۀ ده۔ ځکه که سپوږمۍ له شنې
  پنيرې جوړه وي، نو $2+2=4$۔

  $2+2 = 4$۔ ځکه که سپوږمۍ له شنې پنيرې جوړه وي، نو
  $2+2=4$۔
\end{quote}
خو که «که \dots نو \dots» يوازې $\lif$ وے، جمله به بې له
ستونزې رښتيا، او استنتاجونه به بې له ستونزې معتبر وو۔

بله بېلګه له تاوتولوژۍ $(!A \lif !B) \lor (!B \lif !A)$ سره
تړلې ده۔ دا وړانديز کوي چې که له ورځپاڼې څخه په ناټاکلي ډول
دوه خبري جملې~$S$ او $T$ واخلئ، نو جمله «که $S$، نو $T$؛
يا که $T$، نو~$S$» بايد رښتيا وي۔

\end{document}
